\section{Оценка перспективности научных направлений}

\subsection{Перспективные направления}

\begin{enumerate}
    \item \textbf{Разработка гибридных моделей (MHD + кинетика)} --- сочетание скорости MHD и точности кинетических моделей.
    \item \textbf{Использование полярных сияний для изучения верхней атмосферы} --- спектроскопия сияний позволяет измерять температуру, плотность и состав на высотах 80--300 км.
    \item \textbf{Машинное обучение для прогнозов} --- нейросети, обученные на данных солнечного ветра и геомагнитной активности, могут быстро предсказывать авроральную активность.
\end{enumerate}

\subsection{Направления без большой перспективы}

\begin{enumerate}
    \item \textbf{Поиск единого механизма ускорения} --- данные спутников показывают, что оба механизма (квазистатический и альфвеновский) работают одновременно.
    \item \textbf{Чисто статистические модели без физики} --- они ошибаются при редких выбросах, наиболее опасных для технологий.
    \item \textbf{Увеличение разрешения спутников без улучшения теории} --- важно не количество данных, а их осмысление.
\end{enumerate}